% Slug: tiny-stack-boundary
% Source: chapters/ch01/stack-translation.tex:151 (tikzpicture env start line)
% Original label: \label{fig:tiny-stack-boundary}
% Original caption (LaTeX source):
%   Correct translation connects two execution mechanisms at a common result.
% Original caption (plain text):
%   Correct translation connects two execution mechanisms at a common result.
% Required TikZ setup: \usepackage{tikz} + \usetikzlibrary{arrows.meta, positioning, shapes.geometric, fit} (loaded by support/compiler-pseudocode.sty).
% Shared styles used in this picture: `basic block', `cfg edge'.
%   Defined in support/compiler-pseudocode.sty: `basic block', `cfg edge', `back edge', `block label'.
% Usage: % Slug: tiny-stack-boundary
% Source: chapters/ch01/stack-translation.tex:151 (tikzpicture env start line)
% Original label: \label{fig:tiny-stack-boundary}
% Original caption (LaTeX source):
%   Correct translation connects two execution mechanisms at a common result.
% Original caption (plain text):
%   Correct translation connects two execution mechanisms at a common result.
% Required TikZ setup: \usepackage{tikz} + \usetikzlibrary{arrows.meta, positioning, shapes.geometric, fit} (loaded by support/compiler-pseudocode.sty).
% Shared styles used in this picture: `basic block', `cfg edge'.
%   Defined in support/compiler-pseudocode.sty: `basic block', `cfg edge', `back edge', `block label'.
% Usage: % Slug: tiny-stack-boundary
% Source: chapters/ch01/stack-translation.tex:151 (tikzpicture env start line)
% Original label: \label{fig:tiny-stack-boundary}
% Original caption (LaTeX source):
%   Correct translation connects two execution mechanisms at a common result.
% Original caption (plain text):
%   Correct translation connects two execution mechanisms at a common result.
% Required TikZ setup: \usepackage{tikz} + \usetikzlibrary{arrows.meta, positioning, shapes.geometric, fit} (loaded by support/compiler-pseudocode.sty).
% Shared styles used in this picture: `basic block', `cfg edge'.
%   Defined in support/compiler-pseudocode.sty: `basic block', `cfg edge', `back edge', `block label'.
% Usage: % Slug: tiny-stack-boundary
% Source: chapters/ch01/stack-translation.tex:151 (tikzpicture env start line)
% Original label: \label{fig:tiny-stack-boundary}
% Original caption (LaTeX source):
%   Correct translation connects two execution mechanisms at a common result.
% Original caption (plain text):
%   Correct translation connects two execution mechanisms at a common result.
% Required TikZ setup: \usepackage{tikz} + \usetikzlibrary{arrows.meta, positioning, shapes.geometric, fit} (loaded by support/compiler-pseudocode.sty).
% Shared styles used in this picture: `basic block', `cfg edge'.
%   Defined in support/compiler-pseudocode.sty: `basic block', `cfg edge', `back edge', `block label'.
% Usage: \input{graphs-code/tiny-stack-boundary.tikz.tex} inside a figure environment
%   (wrap with \begin{figure} ... \caption{...} \label{fig:tiny-stack-boundary} ... \end{figure}).
% Note: tikzpicture body below is verbatim from the source; do not reformat.
\begin{tikzpicture}[node distance=15mm and 17mm]
    \node[basic block] (tree) {Expression tree};
    \node[basic block,right=of tree] (code) {Stack code};
    \node[basic block,below=of tree] (value1) {Value in $V$};
    \node[basic block,below=of code] (value2) {Value in $V$};
    \draw[cfg edge] (tree) -- node[above,font=\small] {compile} (code);
    \draw[cfg edge] (tree) -- node[left,font=\small] {evaluate} (value1);
    \draw[cfg edge] (code) -- node[right,font=\small] {execute} (value2);
    \draw (value1) -- node[above,font=\small] {equal} (value2);
  \end{tikzpicture}
 inside a figure environment
%   (wrap with \begin{figure} ... \caption{...} \label{fig:tiny-stack-boundary} ... \end{figure}).
% Note: tikzpicture body below is verbatim from the source; do not reformat.
\begin{tikzpicture}[node distance=15mm and 17mm]
    \node[basic block] (tree) {Expression tree};
    \node[basic block,right=of tree] (code) {Stack code};
    \node[basic block,below=of tree] (value1) {Value in $V$};
    \node[basic block,below=of code] (value2) {Value in $V$};
    \draw[cfg edge] (tree) -- node[above,font=\small] {compile} (code);
    \draw[cfg edge] (tree) -- node[left,font=\small] {evaluate} (value1);
    \draw[cfg edge] (code) -- node[right,font=\small] {execute} (value2);
    \draw (value1) -- node[above,font=\small] {equal} (value2);
  \end{tikzpicture}
 inside a figure environment
%   (wrap with \begin{figure} ... \caption{...} \label{fig:tiny-stack-boundary} ... \end{figure}).
% Note: tikzpicture body below is verbatim from the source; do not reformat.
\begin{tikzpicture}[node distance=15mm and 17mm]
    \node[basic block] (tree) {Expression tree};
    \node[basic block,right=of tree] (code) {Stack code};
    \node[basic block,below=of tree] (value1) {Value in $V$};
    \node[basic block,below=of code] (value2) {Value in $V$};
    \draw[cfg edge] (tree) -- node[above,font=\small] {compile} (code);
    \draw[cfg edge] (tree) -- node[left,font=\small] {evaluate} (value1);
    \draw[cfg edge] (code) -- node[right,font=\small] {execute} (value2);
    \draw (value1) -- node[above,font=\small] {equal} (value2);
  \end{tikzpicture}
 inside a figure environment
%   (wrap with \begin{figure} ... \caption{...} \label{fig:tiny-stack-boundary} ... \end{figure}).
% Note: tikzpicture body below is verbatim from the source; do not reformat.
\begin{tikzpicture}[node distance=15mm and 17mm]
    \node[basic block] (tree) {Expression tree};
    \node[basic block,right=of tree] (code) {Stack code};
    \node[basic block,below=of tree] (value1) {Value in $V$};
    \node[basic block,below=of code] (value2) {Value in $V$};
    \draw[cfg edge] (tree) -- node[above,font=\small] {compile} (code);
    \draw[cfg edge] (tree) -- node[left,font=\small] {evaluate} (value1);
    \draw[cfg edge] (code) -- node[right,font=\small] {execute} (value2);
    \draw (value1) -- node[above,font=\small] {equal} (value2);
  \end{tikzpicture}
